\documentclass[11pt]{article}
        \usepackage[margin=0.72in]{geometry}
        \usepackage[T1]{fontenc}
        \usepackage[utf8]{inputenc}
        \usepackage{lmodern}
        \usepackage{microtype}
        \usepackage{graphicx}
        \usepackage{booktabs}
        \usepackage{tabularx}
        \usepackage{array}
        \usepackage{xcolor}
        \usepackage{float}
        \usepackage{caption}
        \usepackage{enumitem}
        \usepackage{fancyhdr}
        \usepackage[numbers,sort&compress]{natbib}
        \usepackage[colorlinks=true,linkcolor=blue!55!black,citecolor=blue!55!black,urlcolor=blue!55!black]{hyperref}

        \definecolor{navy}{HTML}{163A5F}
        \definecolor{pale}{HTML}{EEF4F9}
        \definecolor{warning}{HTML}{FFF3CD}
        \definecolor{rulegray}{HTML}{B8C2CC}
        \newcommand{\auditbox}[1]{%
          \par\smallskip\noindent
          \fcolorbox{navy}{pale}{\parbox{0.955\linewidth}{\textbf{Detector audit.} #1}}%
          \par\smallskip
        }
        \newcommand{\cautionbox}[1]{%
          \par\smallskip\noindent
          \fcolorbox{rulegray}{warning}{\parbox{0.955\linewidth}{\textbf{Critical limitation.} #1}}%
          \par\smallskip
        }

        \pagestyle{fancy}
        \fancyhf{}
        \lhead{MIT--BIH record 231 beat-type atlas}
        \rhead{RR--HRV front end}
        \cfoot{\thepage}
        \setlength{\headheight}{14pt}
        \setlength{\parindent}{0pt}
        \setlength{\parskip}{4.5pt}
        \setlist[itemize]{leftmargin=1.2em,itemsep=2pt,topsep=3pt}
        \captionsetup{font=small,labelfont=bf}

        \title{\textbf{MIT--BIH Record 231: Detailed Beat-Type Atlas}\\[0.35em]
        \large Expert labels, physiological interpretation, and current QRS-detection audit}
        \author{ECG-only seizure detector project}
        \date{4 August 2026}

        \begin{document}
        \maketitle

        \begin{center}
        \fcolorbox{navy}{pale}{\parbox{0.91\textwidth}{
        \textbf{The central distinction.} The symbols in this report are expert WFDB beat annotations.
        The current RR--HRV front end detects QRS/R events and measures detector agreement; it does
        \emph{not} predict these beat classes. Therefore ``QRS detected'' is never rewritten as
        ``beat type classified correctly,'' and detector disagreement is never automatically called artifact.
        }}
        \end{center}

        \section{Record, data, and selection rule}

        The MIT--BIH Arrhythmia Database contains 48 two-lead ambulatory ECG recordings with expert-reviewed
        beat annotations \cite{goldberger2000physiobank,moody2001mitbih,physionetmitdb}. Record 231 is
        a 30-minute, 360-Hz recording from a 72-year-old woman with MLII and V1. The official notes describe periods of 2:1 AV block, Mobitz II block, rate-related RBBB, and a probable ventricular couplet. The figures use the upper MLII lead.

        The beat labels audited here are \texttt{N} ($n=314$), \texttt{R} ($n=1,254$), \texttt{A} ($n=1$), \texttt{V} ($n=2$). For each label, the displayed target is the chronologically
        middle occurrence after excluding only the 1.5-second file edges. This deterministic rule does not
        select the best-looking or best-detected beat. The comparison comes from another MIT--BIH record with
        the same expert symbol, so the report shows real published variability rather than an invented ideal.

        \subsection{Current detector path and marker meanings}

        The primary event detector is the Khamis/UNSW QRS detector; the NeuroKit gradient detector supplies
        secondary context \cite{khamis2016telehealth,makowski2021neurokit,kristof2024qrs}. The selected
        output is the original-polarity UNSW event refined to the highest positive amplitude within $\pm50$~ms.
        The inverted-polarity result is shown as context only and cannot silently replace the selected timestamp.
        A primary event is called supported when exactly one secondary event lies within $\pm50$~ms. Expert
        matching uses a separate one-to-one $\pm75$-ms audit tolerance.

        In each target panel: the red-edged diamond is the selected expert annotation; gray diamonds and labels
        are neighboring expert beats; the teal triangle is the initial UNSW event; the orange cross is the
        selected refined R timestamp; the purple plus-shaped marker is inverted-path context; and black baseline
        ticks are NeuroKit-gradient context. None of those detector markers is a beat-class prediction.

        \cautionbox{The $\pm75$-ms audit tolerance defines this project's matching experiment; it is not an
        industry declaration that any timestamp inside the window is an exact R-wave fiducial. A wide tolerance
        can hide timing jitter that matters to RR and HRV. Exact raw timing errors are therefore retained.}

        \clearpage
        \section{Full-record QRS-detection audit by expert label}

        \begin{table}[H]
        \centering
        \caption{Selected refined R timestamps versus all expert annotations of each requested label.}
        \small
        \begin{tabularx}{\textwidth}{lXrrrr}
        \toprule
        Symbol & Official meaning & Expert $n$ & Matched/$n$ & Sensitivity & 95\% Wilson interval \\
        \midrule
        \texttt{N} & Expert reference label for a normal beat. & 314 & 314/314 & 100.00\% & 98.79--100.00\% \\
\texttt{R} & Expert reference label for a beat with right bundle-branch-block morphology. & 1,254 & 1,254/1,254 & 100.00\% & 99.69--100.00\% \\
\texttt{A} & Expert reference label for an atrial premature beat. & 1 & 1/1 & 100.00\% & 20.65--100.00\% \\
\texttt{V} & Expert reference label for a premature ventricular contraction. & 2 & 2/2 & 100.00\% & 34.24--100.00\% \\
        \bottomrule
        \end{tabularx}
        \end{table}

        \begin{figure}[H]
        \centering
        \includegraphics[width=0.93\textwidth]{figures/record_231_symbol_sensitivity.png}
        \caption{The interval width, not just the point estimate, exposes how little rare labels establish.
        Intervals use the Wilson score method \cite{wilson1927probable}.}
        \end{figure}

        The full-record initial UNSW result was TP=1,571, FP=11,
        FN=0, and $F_1=0.99651$. After the current positive-amplitude
        fiducial refinement it was TP=1,571, FP=11,
        FN=0, and $F_1=0.99651$. Per-label sensitivity allocates
        matched expert beats to their symbols; false-positive detector events have no expert symbol and therefore
        cannot be represented by sensitivity alone. This table is not a beat-classification confusion matrix.

        \subsection{Exact selected-example audit}
        \begin{table}[H]
        \centering
        \caption{Timing results for the deterministic displayed examples. Errors are detector minus expert.}
        \resizebox{\textwidth}{!}{%
        \begin{tabular}{lrrrrrrrrl}
        \toprule
        Expert & Sample & Time (s) & Initial error (ms) & Selected R error (ms) & Supported? & Inverted R error (ms) & QRS detected? & Beat class output \\
        \midrule
        \texttt{N} & 321,884 & 894.122 & $-2.8$ & $0.0$ & Yes & $-25.0$ & Yes & Not produced \\
\texttt{R} & 380,990 & 1058.306 & $+13.9$ & $0.0$ & Yes & $+47.2$ & Yes & Not produced \\
\texttt{A} & 52,598 & 146.106 & $+11.1$ & $0.0$ & Yes & $+44.4$ & Yes & Not produced \\
\texttt{V} & 53,239 & 147.886 & $+11.1$ & $0.0$ & Yes & $-38.9$ & Yes & Not produced \\
        \bottomrule
        \end{tabular}}
        \end{table}

        \subsection{Record-specific critical reading}
        \begin{itemize}
        \item All requested expert beats were matched, but A has only 1 observation and V only 2. Their 100\% results are case descriptions, not stable class-wide evidence.
\item The refined detector still produced 11 false-positive events despite zero false negatives. Label-wise sensitivity cannot show those errors, which is why full-record PPV/F1 remains necessary.
\item The dominant R labels reflect rate-related RBBB documented by the database. This is a genuine conduction pattern and must not be converted into an artifact label when morphology changes.
        \end{itemize}

        \clearpage
\section{\texttt{N}: normal beat}

\begin{figure}[H]
\centering
\includegraphics[width=\textwidth]{figures/record_231_N_comparison.png}
\caption{Left: the chronologically middle expert \texttt{N} occurrence in record 231
(sample 321,884, 894.122~s), with current detector outputs. Right: an independent
published MIT--BIH record 100 occurrence carrying the same expert symbol. Neighboring
reference symbols are shown to preserve rhythm context \cite{moody2001mitbih,physionetmitdb}.}
\end{figure}

\textbf{Official symbol meaning.} Expert reference label for a normal beat.

\textbf{Why this type occurs and what tends to shape the waveform.} A sinus impulse normally depolarizes the atria first, producing the P wave, then reaches the ventricles through the AV node and His--Purkinje system, producing the QRS; ventricular repolarization produces the T wave \cite{kligfield2007standardization,surawicz2009conduction}. The database symbol describes the expert beat class. It does not require the largest QRS deflection to point upward in every lead.

\textbf{How to read the two strips.} Polarity and amplitude depend on the projection of the cardiac electrical vector onto the selected lead. A negative or biphasic complex may therefore still carry an expert normal-beat label. The comparison strip demonstrates allowed appearance variation; it is not a universal template. The right-hand strip is not a
synthetic textbook ideal and was not selected as a detector success; it is another expert-labeled
observation from the published database. Similar labels can legitimately have different waveforms.

\auditbox{For the deterministic middle N example, the UNSW initial event was $-2.8$~ms from the expert annotation; the selected original-polarity refined timestamp was $0.0$~ms away and was within the predefined $\pm75$-ms audit tolerance. Exact NeuroKit-gradient support within $\pm50$~ms of the primary QRS event was yes. The inverted-context refined timestamp was $-25.0$~ms away. Across the full record, the selected refined output matched 314 of 314 expert N annotations. This is QRS localization, not beat-type classification.}

\cautionbox{A single-lead waveform cannot establish that the entire 12-lead ECG is normal, and QRS detection cannot establish the N class.}
\clearpage
\section{\texttt{R}: right bundle-branch-block beat}

\begin{figure}[H]
\centering
\includegraphics[width=\textwidth]{figures/record_231_R_comparison.png}
\caption{Left: the chronologically middle expert \texttt{R} occurrence in record 231
(sample 380,990, 1058.306~s), with current detector outputs. Right: an independent
published MIT--BIH record 118 occurrence carrying the same expert symbol. Neighboring
reference symbols are shown to preserve rhythm context \cite{moody2001mitbih,physionetmitdb}.}
\end{figure}

\textbf{Official symbol meaning.} Expert reference label for a beat with right bundle-branch-block morphology.

\textbf{Why this type occurs and what tends to shape the waveform.} With RBBB, right-ventricular activation is delayed after left-sided activation. The resulting QRS is prolonged and has lead-dependent terminal changes; formal diagnosis uses a multi-lead pattern \cite{surawicz2009conduction}.

\textbf{How to read the two strips.} The waveform can differ substantially from both N and L beats in the same record. Record 231 is specifically described as having rate-related RBBB, so the conduction pattern may appear or disappear with cycle length. The right-hand strip is not a
synthetic textbook ideal and was not selected as a detector success; it is another expert-labeled
observation from the published database. Similar labels can legitimately have different waveforms.

\auditbox{For the deterministic middle R example, the UNSW initial event was $+13.9$~ms from the expert annotation; the selected original-polarity refined timestamp was $0.0$~ms away and was within the predefined $\pm75$-ms audit tolerance. Exact NeuroKit-gradient support within $\pm50$~ms of the primary QRS event was yes. The inverted-context refined timestamp was $+47.2$~ms away. Across the full record, the selected refined output matched 1254 of 1254 expert R annotations. This is QRS localization, not beat-type classification.}

\cautionbox{The expert R annotation is the reference label. A single MLII strip shows the detector challenge but is not a standalone clinical RBBB diagnosis.}
\clearpage
\section{\texttt{A}: atrial premature beat}

\begin{figure}[H]
\centering
\includegraphics[width=\textwidth]{figures/record_231_A_comparison.png}
\caption{Left: the chronologically middle expert \texttt{A} occurrence in record 231
(sample 52,598, 146.106~s), with current detector outputs. Right: an independent
published MIT--BIH record 100 occurrence carrying the same expert symbol. Neighboring
reference symbols are shown to preserve rhythm context \cite{moody2001mitbih,physionetmitdb}.}
\end{figure}

\textbf{Official symbol meaning.} Expert reference label for an atrial premature beat.

\textbf{Why this type occurs and what tends to shape the waveform.} An ectopic atrial focus fires before the next expected sinus impulse. The early P wave can differ from the sinus P wave or overlap the preceding T wave. If the impulse reaches a recovered His--Purkinje system, the QRS is often narrow and similar to neighboring conducted beats \cite{guichard2022pac,samesima2022ecg}.

\textbf{How to read the two strips.} Prematurity and atrial activity are central to this label; an isolated QRS shape is insufficient. If a bundle branch remains refractory, the same supraventricular impulse can conduct aberrantly and produce a wider or altered QRS. The right-hand strip is not a
synthetic textbook ideal and was not selected as a detector success; it is another expert-labeled
observation from the published database. Similar labels can legitimately have different waveforms.

\auditbox{For the deterministic middle A example, the UNSW initial event was $+11.1$~ms from the expert annotation; the selected original-polarity refined timestamp was $0.0$~ms away and was within the predefined $\pm75$-ms audit tolerance. Exact NeuroKit-gradient support within $\pm50$~ms of the primary QRS event was yes. The inverted-context refined timestamp was $+44.4$~ms away. Across the full record, the selected refined output matched 1 of 1 expert A annotations. This is QRS localization, not beat-type classification.}

\cautionbox{The detector locates ventricular activation. It does not inspect the P-wave origin or prove that an early beat is atrial. The local denominator is only 1; always report the raw 1/1 rather than presenting the percentage alone.}
\clearpage
\section{\texttt{V}: premature ventricular contraction}

\begin{figure}[H]
\centering
\includegraphics[width=\textwidth]{figures/record_231_V_comparison.png}
\caption{Left: the chronologically middle expert \texttt{V} occurrence in record 231
(sample 53,239, 147.886~s), with current detector outputs. Right: an independent
published MIT--BIH record 102 occurrence carrying the same expert symbol. Neighboring
reference symbols are shown to preserve rhythm context \cite{moody2001mitbih,physionetmitdb}.}
\end{figure}

\textbf{Official symbol meaning.} Expert reference label for a premature ventricular contraction.

\textbf{Why this type occurs and what tends to shape the waveform.} A PVC starts in ventricular tissue rather than arriving through the usual atrial/AV conduction sequence. Activation spreads more slowly and asymmetrically through myocardium, commonly creating a broad or unusual QRS with secondary ST--T discordance \cite{prisco2023pvc,surawicz2009conduction}.

\textbf{How to read the two strips.} PVC morphology varies with ventricular origin and lead direction. A broad QRS supports abnormal activation but is not perfectly specific because an atrial or junctional beat may also be wide when conducted aberrantly. The right-hand strip is not a
synthetic textbook ideal and was not selected as a detector success; it is another expert-labeled
observation from the published database. Similar labels can legitimately have different waveforms.

\auditbox{For the deterministic middle V example, the UNSW initial event was $+11.1$~ms from the expert annotation; the selected original-polarity refined timestamp was $0.0$~ms away and was within the predefined $\pm75$-ms audit tolerance. Exact NeuroKit-gradient support within $\pm50$~ms of the primary QRS event was yes. The inverted-context refined timestamp was $-38.9$~ms away. Across the full record, the selected refined output matched 2 of 2 expert V annotations. This is QRS localization, not beat-type classification.}

\cautionbox{The expert V symbol supplies the class. A QRS detector can find the event while remaining incapable of deciding whether its origin was ventricular. The local denominator is only 2; always report the raw 2/2 rather than presenting the percentage alone.}

        \clearpage
        \section{Conclusions for record 231}
        \begin{enumerate}
        \item This report validates QRS-event localization against expert timestamps. It does not validate
        normal, atrial, ventricular, conduction-pattern, or junctional classification;
        seizure detection; or artifact detection.
        \item The full raw denominators and Wilson intervals must accompany percentages. A 100\% result with
        one or two beats remains weak evidence.
        \item The independent comparison strips show that one expert label can have multiple lead-dependent
        shapes. Template resemblance alone is not ground truth.
        \item Original and inverted refinements are both displayed because morphology and polarity can move the
        chosen positive maximum. The inverted path remains context until a routing rule is prospectively specified
        and independently validated.
        \item Detector support is a reliability measurement, not a veto and not an artifact label. Expert
        annotations remain the evaluation reference.
        \end{enumerate}

        \bibliographystyle{unsrtnat}
        \bibliography{MITBIH_Beat_Type_Atlases}
        \end{document}
